\documentclass[10pt,a4paper]{amsart}
\usepackage[margin=19mm]{geometry}
\usepackage{amsmath,amssymb,mathtools}
\usepackage[scr=rsfs,cal=euler]{mathalfa}
\usepackage{xfrac}
\usepackage[unicode,hidelinks,bookmarks=false]{hyperref}
\newcommand{\ie}{i.e.\ }
\newcommand{\cf}{cf.\ }
\newcommand{\resp}{respectively\ }
\newcommand{\Subsection}{\subsection}
\providecommand{\nobreakdash}{\nobreak\hbox{-}}
\setlength{\emergencystretch}{2em}
\multlinegap0pt
\usepackage{amssymb}
\newtheorem{lemma}{Lemma}
\newtheorem{theorem}{Theorem}
\title{Some new estimates for generalized fractional integrals associated with operators on Morrey spaces}
\author{Hua Wang}
\date{Source: arXiv:2605.19372v1 (CC BY 4.0)}
\begin{document}
\maketitle

\noindent\emph{Source and licence.} Some new estimates for generalized fractional integrals associated with operators on Morrey spaces. Version arXiv:2605.19372v1 (CC BY 4.0), distributed by arXiv under the Creative Commons Attribution 4.0 International licence, http://creativecommons.org/licenses/by/4.0/, as recorded in the arXiv licence metadata for this exact version and shown on the abstract page https://arxiv.org/abs/2605.19372v1. Full licence notice: \texttt{licenses/arxiv-2605.19372-CC-BY-4.0.txt}.\par\smallskip

\noindent\textit{Editorial scope.} Counted unit: the article's theorem thm2 (source0.tex lines 493--504). The article's complete original proof of it is delivered with it (source0.tex lines 728--940, one \texttt{proof} environment of 213 lines, which the article places away from the statement). No mathematics is written, repaired or re-derived: the statement and the proof are the article's own lines, in the article's own notation, and no retained line is substituted or re-flowed.  Retained uncounted as the source-local context the proof needs: lines 69--71; lines 135--141; lines 443--446; lines 469--476; lines 485--491; lines 592--598; lines 662--669.  Nothing was omitted from any retained span, so the package carries no omission record.  No retained line contains a citation, so the wrapper carries no bibliography and no bibliography entry.  No retained line contains a figure environment, an image-inclusion command or drawing code, so no image is delivered and the package declares no asset dependency.  Editorial formatting only, in addition: an AMS \texttt{amsart} preamble that restates the article's own declarations and macros used by the retained lines, namely the environments \texttt{lemma}, \texttt{theorem} on separate counters, together with the standard packages those lines need.  The retained environments print in this document as Lemma 1, Theorem 1 in order of appearance, whereas the article prints the same items with the numbering of its own full document.  The complete native source of the article as distributed in the arXiv e-print for this exact version is bundled unchanged as \texttt{sources/200911/source0.tex}.
\begin{equation}\label{G}
\big|\mathcal{P}_t(x,y)\big|\leq\frac{C}{t^{n/2}}\cdot e^{-A\frac{|x-y|^2}{t}}
\end{equation}

\begin{align}\label{kernelk}
\big|\mathcal K_{\alpha}(x,y)\big|
&\leq\frac{1}{\Gamma(\alpha/2)}\int_0^{+\infty}\big|\mathcal{P}_t(x,y)\big|t^{\alpha/2-1}dt\notag\\
&\leq C\cdot\int_0^{+\infty} e^{-A\frac{|x-y|^2}{t}}\cdot t^{\alpha/2-n/2-1}dt\notag\\
&\leq C\cdot\frac{1}{|x-y|^{n-\alpha}}\int_0^{+\infty} e^{-\nu}\cdot\nu^{n/2-\alpha/2-1}d\nu\notag\\
&\leq C\cdot\frac{1}{|x-y|^{n-\alpha}}.
\end{align}

\begin{equation}\label{modulus}
\eta(f;r_{\mathcal{B}}):=\sup_{x_0\in\mathbb R^n}\frac{1}{m(\mathcal{B})}
\int_{\mathcal{B}}\big|f(y)-e^{-t_{\mathcal{B}}\mathcal{L}}f(y)\big|\,dy,
\end{equation}

\begin{lemma}\label{wanglemma2}
Assume that the semigroup $\big\{e^{-t\mathcal L}:t>0\big\}$ has a kernel $\mathcal{P}_t(x,y)$ satisfying the Gaussian upper bound \eqref{G}.
Then for any $0<\alpha<n$, the difference operator $(I-e^{-t\mathcal L})\mathcal{L}^{-\alpha/2}:=\mathcal{L}^{-\alpha/2}-e^{-t\mathcal L}\mathcal{L}^{-\alpha/2}$ has an associated kernel $\widetilde{K}_{\alpha,t}(x,y)$ which satisfies
\begin{equation*}
\widetilde{K}_{\alpha,t}(x,y)\leq\frac{C}{|x-y|^{n-\alpha}}\cdot\frac{t}{|x-y|^2}.
\end{equation*}
Here the constant $C$ is independent of $x,y\in\mathbb R^n$ and $t\in(0,+\infty)$.
\end{lemma}

\begin{theorem}\label{thm1}
Let $0<\alpha< n$ and $1\leq p<n/{\alpha}$. Suppose that $\lambda=n-\alpha p$. Then for any $f\in{M}^{p,\lambda}(\mathbb R^n)$, there exists a positive constant $C>0$ independent of $f$ such that
\begin{equation*}
\big\|\mathcal L^{-\alpha/2}(f)\big\|_{\mathrm{BMO}_{\mathcal{L}}}
\leq C\big\|f\big\|_{{M}^{p,\lambda}}.
\end{equation*}
\end{theorem}

\begin{equation}\label{wang1}
\begin{split}
\int_{\mathcal{B}}\frac{1}{|x-y|^{n-\alpha}}dx&=\int_{|x-y|<3r_{\mathcal{B}}}\frac{1}{|x-y|^{n-\alpha}}dx\\
&=\int_0^{3r_{\mathcal{B}}}\int_{\mathbf{S}^{n-1}}\frac{\varrho^{n-1}}{\varrho^{n-\alpha}}\,d\sigma(x')d\varrho\\
&\leq C(r_{\mathcal{B}})^{\alpha}\leq Cm(\mathcal{B})^{\alpha/n}.
\end{split}
\end{equation}

\begin{equation}\label{wang2}
\begin{split}
\int_{2^{k+1}\mathcal{B}}\frac{1}{|y-z|^{n-\alpha}}dy
&=\int_{|y-z|<2^{k+2}r_{\mathcal{B}}}\frac{1}{|y-z|^{n-\alpha}}dy\\
&=\int_0^{2^{k+2}r_{\mathcal{B}}}\int_{\mathbf{S}^{n-1}}\frac{\varrho^{n-1}}{\varrho^{n-\alpha}}\,d\sigma(y')d\varrho\\
&\leq C\big(2^{k+2}r_{\mathcal{B}}\big)^{\alpha}\leq Cm(2^{k+1}\mathcal{B})^{\alpha/n}.
\end{split}
\end{equation}

\begin{theorem}\label{thm2}
Let $0<\alpha<n$ and $1\leq p<n/{\alpha}$. Suppose that $\lambda=n-\alpha p$. Then for any $f\in V{M}^{p,\lambda}(\mathbb R^n)$, there exists a positive constant $C>0$ independent of $f$ such that
\begin{equation*}
\big\|\mathcal L^{-\alpha/2}(f)\big\|_{\mathrm{VMO}_{\mathcal{L}}}
\leq C\big\|f\big\|_{V{M}^{p,\lambda}},
\end{equation*}
and
\begin{equation*}
\lim_{r_{\mathcal{B}}\to 0^{+}}\eta(\mathcal L^{-\alpha/2}(f);r_{\mathcal{B}})=0.
\end{equation*}
Here $\eta(\mathcal L^{-\alpha/2}(f);r_{\mathcal{B}})$ is defined as in \eqref{modulus}.
\end{theorem}

\begin{proof}[Proof of Theorem $\ref{thm2}$]
Let $f\in V{M}^{p,\lambda}(\mathbb R^n)$ with $1\leq p<n/{\alpha}$ and $\lambda=n-\alpha p$. Since we endow $\mathrm{VMO}_{\mathcal{L}}(\mathbb R^n)$ with the norm of $\mathrm{BMO}_{\mathcal{L}}(\mathbb R^n)$, by Theorem \ref{thm1}, we obtain
\begin{equation*}
\big\|\mathcal L^{-\alpha/2}(f)\big\|_{\mathrm{VMO}_{\mathcal{L}}}
\leq C\big\|f\big\|_{V{M}^{p,\lambda}},
\end{equation*}
where $C>0$ is independent of $f$. We need only to show that
\begin{equation}\label{main22}
\lim_{r_{\mathcal{B}}\to 0^{+}}\eta(\mathcal L^{-\alpha/2}(f);r_{\mathcal{B}})=0.
\end{equation}
Let $\mathcal{B}=B(x_0,r_{\mathcal{B}})$ be a ball centered at $x_0$ and with radius $r_{\mathcal{B}}$. As before, we decompose the function $f(x)$ as follows:
\begin{equation*}
f(x)=f(x)\cdot\chi_{2\mathcal{B}}+f(x)\cdot\chi_{(2\mathcal{B})^{\complement}}:=f_1(x)+f_2(x),
\end{equation*}
where $2\mathcal{B}=B(x_0,2r_{\mathcal{B}})$ and $(2\mathcal{B})^{\complement}=\mathbb R^n\setminus(2\mathcal{B})$. Thus,
\begin{equation*}
\mathcal L^{-\alpha/2}(f)(x)=\mathcal{L}^{-\alpha/2}(f_1)(x)
+\mathcal{L}^{-\alpha/2}(f_2)(x),
\end{equation*}
and
\begin{equation*}
\begin{split}
e^{-t_{\mathcal{B}}\mathcal{L}}\big(\mathcal L^{-\alpha/2}f\big)(x)
&=e^{-t_{\mathcal{B}}\mathcal{L}}\mathcal{L}^{-\alpha/2}(f_1)(x)+
e^{-t_{\mathcal{B}}\mathcal{L}}\mathcal{L}^{-\alpha/2}(f_2)(x),
\end{split}
\end{equation*}
where $t_{\mathcal{B}}=(r_{\mathcal{B}})^2$ and $r_{\mathcal{B}}$ is the radius of the ball $\mathcal{B}$. Then $\eta(\mathcal L^{-\alpha/2}(f);r_{\mathcal{B}})$ can be written as
\begin{equation*}
\begin{split}
&\eta(\mathcal L^{-\alpha/2}(f);r_{\mathcal{B}})\\
&\leq \sup_{x_0\in\mathbb R^n}\frac{1}{m(\mathcal{B})}\int_{\mathcal{B}}\big|\mathcal{L}^{-\alpha/2}(f_1)(x)\big|dx
+\sup_{x_0\in\mathbb R^n}\frac{1}{m(\mathcal{B})}\int_{\mathcal{B}}\big|e^{-t_{\mathcal{B}}\mathcal{L}}\mathcal{L}^{-\alpha/2}(f_1)(x)\big|dx\\
&+\sup_{x_0\in\mathbb R^n}\frac{1}{m(\mathcal{B})}
\int_{\mathcal{B}}\big|\mathcal L^{-\alpha/2}(f_2)(x)-e^{-t_{\mathcal{B}}\mathcal{L}}\mathcal L^{-\alpha/2}(f_2)(x)\big|dx\\
&=J_1(r_{\mathcal{B}})+J_2(r_{\mathcal{B}})+J_3(r_{\mathcal{B}}).
\end{split}
\end{equation*}

\textbf{Estimation of $J_1$:} By using the estimates \eqref{kernelk} and \eqref{wang1}, we have
\begin{equation*}
\begin{split}
J_1(r_{\mathcal{B}})&\leq\sup_{x_0\in\mathbb R^n}\frac{C}{m(\mathcal{B})}\int_{\mathcal{B}}\bigg\{\int_{2\mathcal{B}}\frac{|f(y)|}{|x-y|^{n-\alpha}}dy\bigg\}dx\\
&\leq\sup_{x_0\in\mathbb R^n}\frac{C}{m(\mathcal{B})^{1-\alpha/n}}
\int_{2\mathcal{B}}|f(y)|dy.
\end{split}
\end{equation*}
Applying H\"older's inequality, we obtain
\begin{equation*}
\begin{split}
\frac{1}{m(\mathcal{B})}\int_{2\mathcal{B}}|f(y)|\,dy
&\leq\frac{1}{m(\mathcal{B})}\bigg(\int_{2\mathcal{B}}|f(y)|^p\,dy\bigg)^{1/p}m(2\mathcal{B})^{1/{p'}}\\
&\leq\frac{C}{m(\mathcal{B})^{1/p}}\big\|f\big\|_{L^p(B(x_0,2r_{\mathcal{B}}))}.
\end{split}
\end{equation*}
Since $s\mapsto\|f\|_{L^p(B(x_0,s))}$ is an increasing function and $m(\mathcal{B})^{\alpha/n}\approx(r_{\mathcal{B}})^{\alpha}$, it is easy to see that
\begin{equation*}
\begin{split}
J_1(r_{\mathcal{B}})
&\leq \sup_{x_0\in\mathbb R^n}C(r_{\mathcal{B}})^{\alpha}\int_{2r_{\mathcal{B}}}^{+\infty}\big\|f\big\|_{L^p(B(x_0,s))}\frac{ds}{s^{n/p+1}}\\
&\leq \sup_{x_0\in\mathbb R^n}C(r_{\mathcal{B}})^{\alpha}\int_{2r_{\mathcal{B}}}^{+\infty}\mathcal{M}_{p,\lambda}(f;x_0,s)\cdot m(B(x_0,s))^{\lambda/{pn}}\frac{ds}{s^{n/p+1}}.\\
\end{split}
\end{equation*}
Note that when $\lambda=n-\alpha p$,
\begin{equation}\label{wang5}
\frac{\,n\,}{p}-\frac{\lambda}{\,p\,}=\alpha,
\end{equation}
and $m(B(x_0,s))^{\lambda/{pn}}\approx s^{\lambda/p}$. Therefore,
\begin{equation*}
\begin{split}
J_1(r_{\mathcal{B}})
&\leq C(r_{\mathcal{B}})^{\alpha}\int_{2r_{\mathcal{B}}}^{+\infty}
\sup_{x_0\in\mathbb R^n}\mathcal{M}_{p,\lambda}(f;x_0,s)\frac{ds}{s^{\alpha+1}}.
\end{split}
\end{equation*}
By a change of variables (here we take $\tau:=s/{2r_{\mathcal{B}}}$), we further obtain
\begin{equation*}
J_1(r_{\mathcal{B}})\leq C\int_1^{+\infty}\sup_{x_0\in\mathbb R^n}\mathcal{M}_{p,\lambda}(f;x_0,2r_{\mathcal{B}}\tau)\frac{d\tau}{\tau^{\alpha+1}}.
\end{equation*}
Since
\begin{equation*}
\sup_{x_0\in\mathbb R^n}\mathcal{M}_{p,\lambda}(f;x_0,2r_{\mathcal{B}}\tau)
\leq \sup_{x_0\in \mathbb R^n,r>0}\mathcal{M}_{p,\lambda}(f;x_0,r)=\big\|f\big\|_{{M}^{p,\lambda}}=\big\|f\big\|_{VM^{p,\lambda}},
\end{equation*}
and $\alpha+1>1$, the following integral is convergent:
\begin{equation*}
\int_1^{+\infty}\frac{d\tau}{\tau^{\alpha+1}}<+\infty.
\end{equation*}
Hence,
\begin{equation*}
\lim_{r_{\mathcal{B}}\to 0^{+}}J_1(r_{\mathcal{B}})=0
\end{equation*}
holds by the Lebesgue dominated convergence theorem.

\textbf{Estimation of $J_2$:} From the estimate \eqref{G}, it follows that
\begin{equation*}
\begin{split}
J_2(r_{\mathcal{B}})
&\leq\sup_{x_0\in\mathbb R^n}\frac{1}{m(\mathcal{B})}\int_{\mathcal{B}}
\bigg\{\int_{2\mathcal{B}}\big|\mathcal{P}_{t_{\mathcal{B}}}(x,y)\big|\cdot\big|\mathcal{L}^{-\alpha/2}(f_1)(y)\big|dy\bigg\}dx\\
&+\sup_{x_0\in\mathbb R^n}\frac{1}{m(\mathcal{B})}\int_{\mathcal{B}}
\bigg\{\sum_{k=1}^{\infty}\int_{2^{k+1}\mathcal{B}\setminus 2^k \mathcal{B}}\big|\mathcal{P}_{t_{\mathcal{B}}}(x,y)\big|\cdot\big|\mathcal{L}^{-\alpha/2}(f_1)(y)\big|dy\bigg\}dx\\
&\leq\sup_{x_0\in\mathbb R^n}\frac{C}{m(2\mathcal{B})}\int_{2\mathcal{B}}\big|\mathcal{L}^{-\alpha/2}(f_1)(y)\big|dy\\
&+\sup_{x_0\in\mathbb R^n}
C\sum_{k=1}^{\infty}\frac{1}{2^{kn}}\cdot\frac{1}{m(2^{k+1}\mathcal{B})}\int_{2^{k+1}\mathcal{B}}\big|\mathcal{L}^{-\alpha/2}(f_1)(y)\big|dy\\
&:=J_2'(r_{\mathcal{B}})+J_2''(r_{\mathcal{B}}).
\end{split}
\end{equation*}
Similarly to the estimate of $J_1(r_{\mathcal{B}})$, we have
\begin{equation*}
\lim_{r_{\mathcal{B}}\to 0^{+}}J_2'(r_{\mathcal{B}})=0.
\end{equation*}
On the other hand, from \eqref{kernelk} and \eqref{wang2}, it follows that
\begin{equation*}
\begin{split}
J_2''(r_{\mathcal{B}})&\leq \sup_{x_0\in\mathbb R^n}
C\sum_{k=1}^{\infty}\frac{1}{2^{kn}}\cdot\frac{1}{m(2^{k+1}\mathcal{B})}\int_{2^{k+1}\mathcal{B}}  \bigg\{\int_{2\mathcal{B}}\frac{|f(z)|}{|y-z|^{n-\alpha}}dz\bigg\}dy\\
&\leq \sup_{x_0\in\mathbb R^n}
C\sum_{k=1}^{\infty}\frac{1}{2^{kn}}\cdot\frac{1}{m(2^{k+1}\mathcal{B})^{1-\alpha/n}}\int_{2\mathcal{B}}|f(z)|dz.
\end{split}
\end{equation*}
Moreover, by H\"{o}lder's inequality and the fact that $2\mathcal{B}\subset 2^{k+1}\mathcal{B}$ with $k\in \mathbb{N}$, we get
\begin{equation*}
\begin{split}
\frac{1}{m(2^{k+1}\mathcal{B})}\int_{2\mathcal{B}}|f(y)|\,dy
&\leq\frac{1}{m(2^{k+1}\mathcal{B})}\bigg(\int_{2\mathcal{B}}|f(y)|^p\,dy\bigg)^{1/p}m(2\mathcal{B})^{1/{p'}}\\
&\leq\frac{1}{m(2^{k+1}\mathcal{B})^{1/p}}\big\|f\big\|_{L^p(B(x_0,2r_{\mathcal{B}}))}.
\end{split}
\end{equation*}
Then we have
\begin{equation*}
J_2''(r_{\mathcal{B}})\leq \sup_{x_0\in\mathbb R^n}
C\sum_{k=1}^{\infty}\frac{1}{2^{kn}}\cdot\frac{1}{m(2^{k+1}\mathcal{B})^{1/p-\alpha/n}}\big\|f\big\|_{L^p(B(x_0,2r_{\mathcal{B}}))}.
\end{equation*}
Note that $1/p-\alpha/n>0$. Hence,
\begin{equation*}
\begin{split}
J_2''(r_{\mathcal{B}})&\leq \sup_{x_0\in\mathbb R^n}
C\sum_{k=1}^{\infty}\frac{1}{2^{kn}}\cdot\frac{1}{m(2\mathcal{B})^{1/p-\alpha/n}}\big\|f\big\|_{L^p(B(x_0,2r_{\mathcal{B}}))}\\
&\leq \sup_{x_0\in\mathbb R^n} C(r_{\mathcal{B}})^{\alpha}\int_{2r_{\mathcal{B}}}^{+\infty}\big\|f\big\|_{L^p(B(x_0,s))}\frac{ds}{s^{n/p+1}},
\end{split}
\end{equation*}
since $\|f\|_{L^p(B(x_0,s))}$ is an increasing function with respect to $s$. Arguing as in the proof of $J_1$, we can also obtain
\begin{equation*}
\lim_{r_{\mathcal{B}}\to 0^{+}}J_2''(r_{\mathcal{B}})=0.
\end{equation*}

\textbf{Estimation of $J_3$:} By Lemma \ref{wanglemma2}, we can deduce that
\begin{equation*}
\begin{split}
J_3(r_{\mathcal{B}})&\leq \sup_{x_0\in\mathbb R^n}\frac{1}{m(\mathcal{B})}
\int_{\mathcal{B}}\bigg\{\int_{(2\mathcal{B})^{\complement}}\big|\widetilde{K}_{\alpha,t_{\mathcal{B}}}(x,y)\big|\cdot\big|f(y)\big|dy\bigg\}dx\\
&\leq\sup_{x_0\in\mathbb R^n}\frac{C}{m(\mathcal{B})}
\int_{\mathcal{B}}\bigg\{\sum_{k=1}^{\infty}\int_{2^{k+1}\mathcal{B}\setminus 2^k \mathcal{B}}
\frac{1}{|x-y|^{n-\alpha}}\cdot\frac{(r_{\mathcal{B}})^2}{|x-y|^{2}}\big|f(y)\big|dy\bigg\}dx\\
&\leq \sup_{x_0\in\mathbb R^n}C\sum_{k=1}^{\infty}\frac{1}{2^{2k}}\cdot\frac{1}{m(2^{k+1}\mathcal{B})^{1-\alpha/n}}
\int_{2^{k+1}\mathcal{B}}\big|f(y)\big|dy.
\end{split}
\end{equation*}
Moreover, by H\"{o}lder's inequality, we obtain that for each $k\in \mathbb{N}$,
\begin{equation*}
\begin{split}
\frac{1}{m(2^{k+1}\mathcal{B})}\int_{2^{k+1}\mathcal{B}}|f(y)|\,dy
&\leq\frac{1}{m(2^{k+1}\mathcal{B})}\bigg(\int_{2^{k+1}\mathcal{B}}|f(y)|^p\,dy\bigg)^{1/p}m(2^{k+1}\mathcal{B})^{1/{p'}}\\
&=\frac{1}{m(2^{k+1}\mathcal{B})^{1/p}}\big\|f\big\|_{L^p(B(x_0,2^{k+1}r_{\mathcal{B}}))}.
\end{split}
\end{equation*}
From this, it follows that
\begin{equation*}
J_3(r_{\mathcal{B}})\leq\sup_{x_0\in\mathbb R^n}C\sum_{k=1}^{\infty}\frac{1}{2^{2k}}\cdot\frac{1}{m(2^{k+1}\mathcal{B})^{1/p-\alpha/n}}
\big\|f\big\|_{L^p(B(x_0,2^{k+1}r_{\mathcal{B}}))}.
\end{equation*}
Since $s\mapsto\|f\|_{L^p(B(x_0,s))}$ is an increasing function, we have
\begin{equation}\label{wang6}
\begin{split}
\int_{4r_{\mathcal{B}}}^{+\infty}\frac{\|f\|_{L^p(B(x_0,s))}}{s^{n/p-\alpha+3}}ds
&=\sum_{k=1}^{\infty}\int_{2^{k+1}r_{\mathcal{B}}}^{2^{k+2}r_{\mathcal{B}}}\frac{\|f\|_{L^p(B(x_0,s))}}{s^{n/p-\alpha+3}}ds\\
&\geq\sum_{k=1}^{\infty}\int_{2^{k+1}r_{\mathcal{B}}}^{2^{k+2}r_{\mathcal{B}}}
\big\|f\big\|_{L^p(B(x_0,2^{k+1}r_{\mathcal{B}}))}\frac{ds}{s^{n/p-\alpha+3}}\\
&\geq C\sum_{k=1}^{\infty}\big\|f\big\|_{L^p(B(x_0,2^{k+1}r_{\mathcal{B}}))}\cdot\frac{1}{(2^{k+1}r_{\mathcal{B}})^{n/p-\alpha+2}}\\
&\geq \frac{C}{(r_{\mathcal{B}})^2}\sum_{k=1}^{\infty}\big\|f\big\|_{L^p(B(x_0,2^{k+1}r_{\mathcal{B}}))}\cdot
\frac{1}{2^{2k}}\cdot\frac{1}{m(2^{k+1}\mathcal{B})^{1/p-\alpha/n}}.
\end{split}
\end{equation}
Therefore, by \eqref{wang5}, \eqref{wang6} and the definition of Morrey norm, we deduce that
\begin{align*}
J_3(r_{\mathcal{B}})&\leq \sup_{x_0\in\mathbb R^n}
C(r_{\mathcal{B}})^2\int_{4r_{\mathcal{B}}}^{+\infty}\frac{\|f\|_{L^p(B(x_0,s))}}{s^{n/p-\alpha+3}}ds\\
&\leq \sup_{x_0\in\mathbb R^n}
C(r_{\mathcal{B}})^2\int_{4r_{\mathcal{B}}}^{+\infty}\frac{\mathcal{M}_{p,\lambda}(f;x_0,s)\cdot s^{{\lambda}/p}}{s^{n/p-\alpha+3}}ds\\
&\leq C(r_{\mathcal{B}})^2\int_{4r_{\mathcal{B}}}^{+\infty}\sup_{x_0\in\mathbb R^n}\mathcal{M}_{p,\lambda}(f;x_0,s)\frac{ds}{s^3}.
\end{align*}
Observe that by a change of variables (here we take $\tau:=s/{4r_{\mathcal{B}}}$),
\begin{equation*}
(r_{\mathcal{B}})^2\int_{4r_{\mathcal{B}}}^{+\infty}\sup_{x_0\in\mathbb R^n}\mathcal{M}_{p,\lambda}(f;x_0,s)\frac{ds}{s^3}
=\int_1^{+\infty}\sup_{x_0\in\mathbb R^n}\mathcal{M}_{p,\lambda}(f;x_0,4r_{\mathcal{B}}\tau)\frac{d\tau}{\tau^3}.
\end{equation*}
Since for any $r_{\mathcal{B}}>0$,
\begin{equation*}
\sup_{x_0\in\mathbb R^n}\mathcal{M}_{p,\lambda}(f;x_0,4r_{\mathcal{B}}\tau)
\leq\sup_{x_0\in \mathbb R^n,r>0}\mathcal{M}_{p,\lambda}(f;x_0,r)=\big\|f\big\|_{{M}^{p,\lambda}}
=\big\|f\big\|_{VM^{p,\lambda}},
\end{equation*}
and the following integral is convergent:
\begin{equation*}
\int_1^{+\infty}\frac{d\tau}{\tau^3}<+\infty.
\end{equation*}
Finally, by using the Lebesgue dominated convergence theorem, we conclude that
\begin{equation*}
\lim_{r_{\mathcal{B}}\to 0^{+}}J_3(r_{\mathcal{B}})=0.
\end{equation*}
Combining the estimates of $J_1$, $J_2$ with $J_3$, we get the desired result \eqref{main22}. Hence, the proof of Theorem \ref{thm2} is complete.
\end{proof}

\end{document}
